\chapter{伯努利流形上的可解动力学：坐标、边界与目标势}
\label{chp:analytical_solution_bernoulli}

本章给出一个可以逐项复核的低维基准，而不是把一维模型包装成整套智能理论的证明。我们把二元信念族写成带 Fisher 度量的状态空间，依次定义状态、时间、代价单位、作用量、边界条件和观测接口，再讨论连续动力学、离散谱以及目标势。这个基准保留了 HSF-HD 所关心的状态变化、结构约束、外部输入和读出，却有意舍弃真实智能系统的大部分高维耦合。它的价值在于暴露量纲错误、边界偷换和过度本体化：解析可解只能说明某组假设下的模型性质，不能由此推出神经系统、语言模型或 AgentOS 与该模型整体同构。AgentOS 中的 $h(t)/E/\Theta$、$\Psi$、SPC、TCE、CCM、Boundary、Offloading 与外部记忆 $M_e$ 可用它做单元测试，但它们只是工程实现之一。

\section{从二元信念到 Fisher 坐标}
\label{sec:dimensional_reduction_bernoulli}

先定义对象。观测变量 $X\in\{0,1\}$，状态参数 $q\in(0,1)$，分布为 $p(X=x\mid q)=q^x(1-q)^{1-x}$。这里的 $q$ 是对一个已经固定身份、固定语义和固定时间窗的二元命题所作的概率读出，不是智能体全部内部状态。命题身份必须由任务接口绑定，例如“传感器在未来两秒内是否越过阈值”；若身份、阈值或时间窗随时变化，就应把它们并入扩展状态或外部输入，而不能仍把变化归因于单一的 $q(t)$。局部结构网络保存实体关系和约束，全局分布表示提供上下文相关的连续编码；二者经身份绑定、耦合与读出接口共同产生这个标量。把二者压成一个概率是模型约化，不意味着某个节点、脑区或参数就是该信念本身。

对伯努利族，得分函数为 $\partial_q\log p(X\mid q)=X/q-(1-X)/(1-q)$。在样本由当前分布产生、且支持集不随 $q$ 改变的假设下，Fisher 信息为
\[
g(q)=\mathbb E_q[(\partial_q\log p)^2]=\frac{1}{q(1-q)}.
\]
因此无量纲线元为 $ds^2=g(q)dq^2$。令
\begin{equation}\label{eq:bernoulli_transform}
q=\sin^2\!\left(\frac{\theta}{2}\right),\qquad \theta\in(0,\pi),
\end{equation}
则 $dq=\tfrac12\sin\theta,d\theta$，而 $q(1-q)=\tfrac14\sin^2\theta$，于是严格得到 $ds^2=d\theta^2$。这一步也纠正一个常见误读：虽然 $g(q)$ 在 $q\to0,1$ 时发散，端点到任意内点的测地距离却有限。例如从 $q=1/2$ 到 $q=1$ 的距离是 $\int_{1/2}^{1}dq/\sqrt{q(1-q)}=\pi/2$，所以“度量分量发散”不能推出“绝对信念位于无限远处”。端点难以改变若成立，必须来自控制饱和、数据稀缺、先验、边界策略或结构耦合，而不是这一坐标奇性本身。

坐标 $q$ 与 $\theta$ 描述同一统计状态，任何物理或计算结论都应在重参数化后保持一致。比如速度的数值会改变：$\dot q=(\sin\theta/2)\dot\theta$，但路径长度 $\int\sqrt{g(q)}|\dot q|dt=\int|\dot\theta|dt$ 不变。工程日志若只记录 $|\dot q|$，会在端点附近误报“停滞”；记录 Fisher 速度 $v_F=\sqrt{g(q)}|\dot q|$ 才能跨坐标比较。反过来，Fisher 距离只度量分布可区分性，不等于任务损失、信息熵、活动强度或焦耳能耗。若错误答案以高置信度给出，几何路径可以很短而任务失败仍然严重。这个反例提醒我们：状态度量回答“分布改变多少”，目标函数回答“改变是否有用”，两者必须分别定义。

把该约化用于数据时还要注明估计层。有限样本给出的 $\hat q=k/n$ 可能落在端点，使经验 Fisher 量数值爆炸；常用的 Beta 平滑、截断 $q\in[\varepsilon,1-\varepsilon]$ 或自然参数 $\eta=\log(q/(1-q))$ 都会改变估计偏差。它们是数值设计，不是理论定理。检验时应报告样本量、平滑参数、置信区间和命题漂移，并以校准误差或对数损失验证概率读出。只有身份稳定、观测模型合适且估计误差可控，后续一维动力学才具有可解释对象。

这一约化还要求区分总体参数、主体信念与观测频率。若 $q$ 是环境事件的总体发生率，更新由采样模型决定；若 $q$ 是智能体的主观置信度，还需给出从内部状态到概率报告的校准接口；若 $q$ 只是有限窗口中的经验频率，它会随窗口滚动而机械变化。三者数值相同也不代表因果机制相同。一个可靠实验应并行保存原始观测、模型后验和最终报告，并通过干预读出器或替换先验来区分来源。对于多主体任务，还要声明每个主体是否共享同一命题身份和证据集；未经对齐的概率不能直接平均。若命题从“会下雨”悄然改成“本区域会下雨”，状态空间已改变，连续轨迹应在事件处重置或通过显式映射衔接。

从高维状态投影到 $q$ 也会遗失信息。设原状态为 $x\in\mathcal X$，读出为 $q=R_\iota(x)$，则不同 $x$ 可以给出同一 $q$；只有当未来任务转移和损失在这些等价状态上近似相同，$q$ 才是任务充分统计量。检验方法是从相同 $q$、不同结构记录的样本出发，比较后续响应分布。若响应显著不同，就必须扩展状态，而不能通过调节一维势函数掩盖缺失变量。这个充分性检查把“为了可解而降维”与“系统真的只有一维”明确分开。

\section{作用量、惯性参数与量纲闭合}
\label{sec:cognitive_inertia_nert_mass}

为避免把数学符号直接冒充物理量，先选定工程单位。令 $t$ 以秒计，$q$ 与 $\theta$ 无量纲；令 $C$ 表示一次任务定义下的“计算代价单位”，它可以由延迟、调用费、写入次数和风险罚项经明确权重组合，但不等同于纳特、焦耳或香农熵。定义路径作用量
\[
S[q]=\int_{t_0}^{t_1}\left[\frac{\mu}{2}g(q)\dot q^2-V(q,t)\right]dt .
\]
若 $S$ 的单位取 $C$，则拉格朗日量单位为 $C/\mathrm{s}$，$g(q)\dot q^2$ 的单位为 $\mathrm{s}^{-2}$，所以惯性参数 $\mu$ 的单位是 $C\!\cdot\!\mathrm{s}$；势函数 $V$ 的单位是 $C/\mathrm{s}$。这套量纲闭合足以支持比较和实现，没有必要把 $\mu$ 宣称成新的自然基本单位，也不能从形式计算推出它等于时间、网络大小或物质质量。

在 $V=0$ 且结构、度量不显含时间时，共轭量为 $p_q=\mu g(q)\dot q$，哈密顿函数为 $H=p_q^2/[2\mu g(q)]$。换到 $\theta$ 坐标得到 $p_\theta=\mu\dot\theta$ 与 $H=p_\theta^2/(2\mu)$。两种写法一致，说明 $q(1-q)$ 出现在 $q$ 坐标的哈密顿式中只是逆度量因子；不能据其在端点趋零就断言“能量冰封”。保持固定 $p_q$ 与保持固定几何速度是不同实验条件，前者甚至会让坐标速度随位置改变。任何端点结论都必须说明固定了什么、由何种控制器施力以及边界是否可达。

欧拉—拉格朗日方程也应连续推导。若 $V=V(q,t)$，则
\[
\mu g(q)\ddot q+\frac{\mu}{2}g'(q)\dot q^2+\partial_qV(q,t)=0.
\]
若目标、结构或输入变化，导数中不能漏项。更一般地，$g=g(q,z(t),t)$ 且有外部控制 $u(t)$ 时，方程还包含 $\mu\partial_tg\,\dot q$、$\mu(\partial_zg\cdot\dot z)\dot q$ 和控制广义力 $B(q,t)u(t)$。只写静态梯度下降会把环境漂移错误归因于内部意图。在 $\theta$ 坐标中同一系统写作 $\mu\ddot\theta+\partial_\theta\widetilde V=\widetilde B u$，可作为实现的交叉检查。

离散实现必须声明步长。以半隐式更新为例，
\[
v_{k+1}=v_k-\Delta t\,\mu^{-1}\partial_\theta\widetilde V(\theta_k,t_k),\qquad
\theta_{k+1}=\theta_k+\Delta t\,v_{k+1}.
\]
稳定性依赖势曲率、阻尼和 $\Delta t$，不能从连续模型自动继承。若增加阻尼 $\gamma v$，则 $\gamma$ 的单位为 $C$，代价率沿轨迹减少的是所定义的机械型函数，而非必然减少实际电耗或任务错误。一个控制器可以持续降低 $V$，却因目标设错而更快走向错误答案；也可能降低代理损失但破坏校准。故验收至少同时报告几何轨迹、任务成功、资源账本与边界违规，不能用单一“能量下降”替代成功证明。

在 AgentOS 中，$\mu$ 可作为经实验辨识的有效参数：对固定任务施加已知幅度的提示、检索或工具反馈，观测 $\theta(t)$ 的响应，再拟合动力学。它会随上下文长度、外部记忆 $M_e$、模型版本和工具延迟而变，是条件参数而非智能体本体属性。可辨识性检验要比较不同 $\mu,V,B$ 组合是否产生近似相同轨迹；若存在等价组合，就只能报告参数区间。这样，“认知惯性”保留了工程意义，同时不再借量纲修辞制造虚假的自然常数。

参数辨识还应处理采样率与未观测输入。若日志每十秒记录一次，而控制在秒级变化，二阶差分会把高频输入折叠成表观惯性；若模型调用排队时间未入账，拟合出的 $\mu$ 会吸收调度延迟。可采用不同采样率、已知脉冲输入和保持实验作交叉验证，并将过程噪声与测量噪声分开。报告拟合优度之外，还应给出残差自相关、参数置信区间和跨任务外推误差。只有响应随输入幅度、初值和时间尺度呈现一致规律时，惯性参数才是有用简化。

资源账本必须独立于作用量账本。一次检索可增加延迟却降低错误风险，一次长推理可提高令牌消耗却减少后续回滚，设备焦耳能耗又受硬件和并发影响。工程代价 $C$ 可以通过决策者声明的权重把这些量归一化，但权重变化会改变 $\mu$ 与 $V$ 的数值，因此需要版本控制。若目标是比较两台设备的实际能效，就应直接测量功率积分；若目标是比较策略质量，就用任务代价。两种结果可以相关，不能互相替代。

\section{结构耦合、编辑传播与有限延迟}
\label{sec:cognitive_equivalence_principle}

一维状态并不取消底层结构。设局部结构网络为 $G_s=(\mathcal V,\mathcal E,W)$，节点表示任务中可寻址的实体或约束，边表示可审计的关系；设全局分布表示为 $h\in\mathbb R^d$，它编码上下文相关但不必逐节点可解释的状态。绑定算子 $A_\iota$ 根据任务身份 $\iota$ 从二者取得证据，读出函数 $R_\iota(G_s,h)$ 产生 $q$。因而 $q$ 的改变可能来自节点编辑、边权改变、全局表示移动、读出器更新或输入改变。没有干预实验时，这些机制不可由同一条 $q(t)$ 曲线区分，更不能从“响应慢”直接推出网络规模大。

为讨论传播下界，给每条有向边 $e$ 配置非负传输与处理延迟 $\tau_e$，并规定一次编辑必须使受影响集合 $U\subseteq\mathcal V$ 的约束达到一致。若输入首先抵达节点 $s$，任何遵守因果消息传递的实现都满足
\[
T_{\rm settle}\ge \max_{v\in U}\min_{P:s\leadsto v}\sum_{e\in P}\tau_e .
\]
这是图上最短延迟路径的最大值，而不是含糊的“拓扑体积除以认知光速”。它的假设可测试：消息不能跳过接口、边延迟有下界、所有 $U$ 中节点都必须确认。若系统允许广播、共享内存、并行参数更新或直接重算全局表示，下界会改变；若只需局部读出正确，也不必等待全图收敛。因此传播时间与结构范围相关，却不与惯性参数 $\mu$ 绝对同一。

以“把某知识库中的首都关系从甲改为乙”为例，局部关系表可在一个事务内修改，但全局向量索引、缓存、派生问答、权限规则和外部记忆仍可能保留旧值。安全过程应先冻结任务身份，列出依赖闭包，再分阶段执行写入、重嵌入、缓存失效和一致性探测。探测不能只问原句，还应覆盖逆向关系、组合问题、否定问题和无关事实。若主问题已答对而无关事实大幅退化，任务成功仍不成立。这里所谓“拔出萝卜带出泥”被还原成依赖传播与回归风险，而不是未经定义的拓扑撕裂。

对应到 AgentOS，SPC 可以负责提出候选状态，TCE 管理目标与期限，CCM 维护一致性检查，Boundary 限制可写范围，Offloading 把长链证据交给外部记忆 $M_e$。这些模块共同给出一种实现，但 HSF-HD 的一般状态动力学只要求状态、结构、输入、约束和读出被显式区分，不要求所有智能系统拥有同名模块。尤其不能把局部结构网络等同于全局分布表示：前者适合事务一致性和因果追踪，后者适合相似性与上下文泛化。二者通过身份键、时间戳、版本向量和读出接口联合，错误绑定会造成比数值误差更严重的对象错置。

工程上可把传播试验做成阶跃响应。记录编辑开始、各依赖节点确认、读出首次正确、全套回归通过四个时间点，并在不同图直径、并行度和缓存策略下重复。若延迟随最远依赖路径增加，支持有限传播模型；若延迟主要由批处理周期决定，则应修改机制解释。还应设置反例：一个大而稀疏但可广播的图可能比小而串行的图更快，一个强耦合表示也可能通过低秩更新快速改变。只有在控制混杂变量后，才能把有效惯性归因于结构。这样保留了“复杂信念需要协调”的洞见，同时拒绝时间、质量和拓扑大小的无条件等号。

上述下界也不等于完成时间的预测值。真实完成时间还包含排队、失败重试、权限审批、批处理栅栏和人工复核，而某些依赖可并行取消。应把因果传播下界、调度上界和实测分布分别报告；否则一次网络拥塞就会被误解为概念更“沉重”。这一区分让结构机制能够被证伪，也为缓存和广播优化留下明确位置。

\section{边界条件、谱与离散模式}
\label{sec:topological_boundaries_and_quantization}

解析谱首先是算子与定义域的共同性质。利用 $\theta\in(0,\pi)$ 的等距坐标，Hilbert 空间取 $\mathcal H=L^2((0,\pi),d\theta)$，内积为 $\langle f,g\rangle=\int_0^\pi f^*g,d\theta$。定义非负算子 $L=-d^2/d\theta^2$，但在指定边界前它还不是完整的谱问题。若选择 Dirichlet 条件 $f(0)=f(\pi)=0$，定义域可取 $H^2(0,\pi)\cap H_0^1(0,\pi)$。积分分部的边界项消失，$L$ 自伴，其归一化本征函数与本征值为
\[
\phi_n(\theta)=\sqrt{2/\pi}\sin(n\theta),\qquad \lambda_n=n^2,quad n=1,2,\ldots .
\]
换回 $q$ 即得到 $\phi_n(q)=\sqrt{2/\pi}\sin[2n\arcsin\sqrt q]$。该离散性来自有限区间、算子和边界条件的组合。

边界不是由“绝对真假”自动强迫出来的。若取 Neumann 条件 $f'(0)=f'(\pi)=0$，谱包含常数零模与 $\cos(n\theta)$；若把两端周期识别，则得到复指数模并出现不同简并；Robin 条件还引入连续可调的边界参数。若端点允许概率质量吸收、反射或重置，相应生成元又不同。这些反例说明 Dirichlet 正弦模不是唯一自然答案，也不能被称为“绝对语义子”。真实 token 的离散性主要来自编码器、词表与采样接口；连续表示中的谱模可以帮助分析，却不构成词元存在的本体证明。

若用二阶扩散或波动模型，系数仍需定义。例如 $\partial_tu=-\kappa Lu$ 中 $\kappa$ 单位为 $\mathrm{s}^{-1}$，模态按 $e^{-\kappa n^2t}$ 衰减；若写 $\partial_t^2u=-c^2Lu$，则 $c$ 单位为 $\mathrm{s}^{-1}$，频率为 $cn$。若进一步定义工程哈密顿算子 $\widehat H=(\alpha/2)L$，其本征“代价率”为 $\alpha n^2/2$，只有当 $\alpha$ 的单位和解释明确时才有意义。把 $\alpha$ 写成类似 $\hbar^2/\mu$ 的形式只是参数化选择，不能据此宣称发生物理量子化。

谱基底仍然很有用。任意满足相应边界和正则性的状态可展开为 $u(\theta,t)=\sum_na_n(t)\phi_n(\theta)$，低阶模表示平滑的大尺度变化，高阶模表示快速振荡。压缩、滤波和稳定性分析可以在模态系数上完成。可是模态编号不等于语义复杂度：高阶振荡可能只是噪声，低阶模也可能承载关键决策。若读出 $R$ 只观察某些线性组合，未被读出的模仍会影响后续动力学。故工程报告需同时给出谱能量、任务读出和残差，而不能把 $n$ 直接命名为“概念等级”。

数值验证可用三组边界并行进行。将区间离散为 $N$ 个网格，声明 $\Delta\theta=\pi/(N+1)$ 或相应端点方案，计算矩阵本征值并与解析谱比较；随后做网格加密，确认误差按离散阶收敛。再把同一初态置于 Dirichlet、Neumann 和 Robin 条件下，比较长期行为。若软件在三者间给出同一谱，说明边界实现被忽略。最后检查离散算子的加权对称性和概率/范数预算。可解模型的价值正是在这里：它提供可重复的基准，而非宏大结论的捷径。

在 $q$ 坐标直接离散时尤其容易出错。均匀 $q$ 网格不是均匀 Fisher 网格，端点附近的系数变化会造成条件数恶化；若简单套用中心差分，离散矩阵可能不再相对于正确权重对称。稳妥做法是先在等距 $\theta$ 坐标离散，或用有限体积形式保留通量和测度，再把结果映回 $q$。两种实现应在网格加密后收敛到同一低阶谱。若研究吸收边界，则还需监测总质量流出并把它记入预算，不能把范数下降解释为内部耗散。

谱截断也需要误差界。保留前 $K$ 个模态时，截断残差由被舍弃系数决定，且目标读出可能对某些高频模异常敏感。因此不能只依据累计谱能量选 $K$，还要测量任务损失和边界量的误差。对非光滑初态会出现 Gibbs 振荡，可通过弱解、滤波或更合适的基底处理，但每种处理都会改变短时行为。工程结论应注明保留模态数、滤波器与误差容限，使“低维有效”成为可复核的近似而非视觉判断。

\section{一阶生成元、自伴定义域与可观测性}
\label{sec:non_self_adjoint_momentum_blind_spot}

考虑形式上的一阶算子 $P=-i\alpha\,d/d\theta$，其中 $\alpha$ 的单位由应用决定。仅写这个微分表达式无法判断它是否为可观测量；必须给出 Hilbert 空间与定义域。对足够光滑的 $f,g$，
\[
\langle f,Pg\rangle-\langle Pf,g\rangle=-i\alpha[f^*(\theta)g(\theta)]_0^\pi .
\]
在两端都取零的 Dirichlet 定义域上，$P$ 虽对称，却通常不是自伴，因为伴随域更大；而满足 $f(\pi)=e^{i\beta}f(0)$ 的准周期定义域可给出一族自伴扩张。二阶算子 $P^2$ 在 Dirichlet 条件下可以自伴，并不要求 $P$ 在同一域上也自伴。由此可见，所谓“动量非自伴”是边界与定义域的具体数学结果，不是关于意识层级的普遍事实。

数学上的自伴性与统计上的可观测性也必须分开。自伴性保证理想封闭模型中的实谱和酉生成性质；可观测性则问给定输入、传感器和噪声，内部状态能否从输出恢复。对线性系统 $\dot x=Ax+Bu,\ y=Cx+\epsilon$，有限维可观测性可由 $[C;CA;\ldots;CA^{n-1}]$ 的秩判断。即使某算子自伴，若 $C$ 抹去相关方向，状态仍不可辨识；反之，非正规或开放系统的参数也可能从时间序列稳定估计。把前者直接叫作“反思盲区”会把算子域问题偷换成心理经验主张。

在 HSF-HD 语境中，盲区更适合定义成接口性质。设完整状态为 $x=(h,G_s,M_e,z)$，可见读出为 $y=R(x)$。若存在不同状态 $x_1\ne x_2$ 在允许的全部探测序列下产生不可区分输出，则这些状态位于当前接口的不可辨识等价类。增加工具、记忆检索、反事实提示或结构审计可能细分该等价类；因此盲区会随实验协议改变。它不是一个永远封闭的“潜意识区域”，也不能靠脑区类比证明。局部结构与全局分布的联合探测尤其重要：只读自然语言报告可能漏掉结构冲突，只查图数据库又可能漏掉上下文漂移。

工程上应实施输入设计。对候选模式施加不同频率和方向的小扰动，记录输出灵敏度矩阵；用留出轨迹检验重建器，而不是在拟合数据上宣布可观测。若两个模型都解释观测，应报告模型集合或后验区间。还要防止探测改变被测状态：长链追问、工具调用和记忆写回都会成为输入，故测量协议应进入动力学方程。若目标 $C(t)$、读出 $R_t$ 或结构 $G_s(t)$ 在实验中变化，导数和辨识模型都要包含这些变化，不能把非平稳性归咎于隐藏意志。

一个具体反例是：系统内部保留了正确证据，但生成接口因安全边界拒绝输出。只看文本会判断“知识不存在”，检查授权后的结构读出却可恢复证据。另一个反例是自然语言报告正确，而依赖图仍含冲突旧边，后续组合任务因而失败。这两种情形说明可观察输出、内部一致性与任务成功是三个不同命题。用伯努利基准分析一阶生成元，可以精确揭示边界项；要解释智能系统的反思能力，则还需要实验接口、噪声模型和因果干预。

对开放系统，一阶生成元还可能含漂移、扩散、跳变和吸收项。此时概率密度的演化由主方程或 Fokker--Planck 型算子描述，适当问题通常是守恒性、正性与稳态，而不是寻找一个封闭系统的“动量观测量”。离散工具调用会导致状态跳变，权限拒绝会导致边界反射或终止，记忆写入会改变后续生成元。把这些事件硬塞进固定的一阶微分算子，会产生看似精确却对象错误的谱。混合系统应分别定义连续段、事件守卫和重置映射，再检验整体可观测性。

可观测性也具有尺度依赖。短时间窗内两个慢变量可能完全不可区分，延长实验后才显现；反之，环境漂移会让过长窗口失去平稳性。应在多个时域上计算经验可观测 Gramian 或信息矩阵，并用奇异值而非二元标签描述弱方向。若最小奇异值接近噪声底，只能说当前协议辨识能力弱，不能宣布相关状态原则上不可知。

\section{目标势、局部近似与 Mathieu 方程}
\label{sec:teleological_potential_quantum_pendulum}

目标导向动力学可以通过势函数表达，但势必须由任务定义。设目标概率 $q_*\in(0,1)$ 固定，选取
\[
V(q)=\rho\,D_{\mathrm{KL}}(\mathrm{Bern}(q_*)\Vert\mathrm{Bern}(q)),
\]
其中 $\rho$ 的单位为 $C/\mathrm{s}$，KL 散度无量纲。这个方向的 KL 在 $q$ 错误逼近端点时发散，体现概率校准罚项；交换 KL 方向会得到不同曲率与边界行为，不能混用。若目标 $q_*=q_*(t)$ 变化，则 $V(q,t)$ 显含时间，沿轨迹的总导数含 $\partial_tV$，即使 $q$ 沿负梯度移动，$V$ 也可能因目标迁移而上升。

取对称目标 $q_*=1/2$，代入 $q=\sin^2(\theta/2)$ 可得
\[
D_{\mathrm{KL}}(1/2\Vert q)=-\log(\sin\theta)+\text{常数}.
\]
这不是余弦势，也不直接给出 Mathieu 方程。只有在指定工作点附近做局部 Taylor 近似，或另行选择周期目标势 $\widetilde V(\theta)=V_0[1-\cos(2\theta)]/2$ 时，线性谱方程才可化为标准 Mathieu 形式。若考虑
\[
\left[-\frac{\alpha}{2}\frac{d^2}{d\theta^2}+\frac{V_0}{2}(1-\cos2\theta)\right]\psi=\varepsilon\psi ,
\]
整理后得到 $\psi''+[a-2r\cos(2\theta)]\psi=0$，其中 $a=(2\varepsilon-V_0)/\alpha$、$r=-V_0/(2\alpha)$。参数映射、区间与边界必须同时给出，才能调用 Mathieu 特征值；这是一类条件模型，不是“意志必然形成量子单摆”。

若采用经典受控轨迹而非谱方程，则在 $\theta$ 坐标写成
\[
\mu\ddot\theta+\gamma\dot\theta+\partial_\theta\widetilde V(\theta,t)=b(\theta,t)u(t)+\xi(t).
\]
$\xi$ 是需声明统计性质的扰动，不能自动解释为量子涨落。局部稳定性由平衡点处势曲率与阻尼决定；势下降只证明代理函数改善，不证明任务完成。任务成功还要求读出越过验收阈值、约束未违反、校准可接受，并在目标变化后保持鲁棒。若控制器把 $q$ 推到目标值却同步损坏结构网络中的身份绑定，表面轨迹正确而系统已经回答了另一个问题。

目标势还可以表达多目标权衡。令 $V=\rho_1L_{\rm task}+\rho_2L_{\rm consistency}+\rho_3L_{\rm risk}$，各项先无量纲化后再组合，权重才有可解释性。若直接把延迟毫秒、KL 纳特、货币成本和安全等级相加，数值下降没有稳定含义。Boundary 可把不可接受区域写成障碍或投影，Offloading 可通过外部记忆减少内部轨迹长度，但二者会改变可行域和输入，应进入模型。TCE 调整目标时必须留下版本与时间戳，使 $\partial_tV$ 可追踪；CCM 则对结构一致性作独立验收。这里的模块名服务于 AgentOS 实例，不构成一般智能的必要器官表。

验证应比较精确 KL 势、局部二次近似和周期余弦势。选取多个初值，计算轨迹与谱的差异，报告近似在何种 $|\theta-\theta_*|$ 范围内误差低于阈值。再加入目标突变、输入延迟与噪声，检查控制器是否把非平稳项正确计入。若 Mathieu 近似只在狭小区域有效，就把它标记为局部求解器；若边界附近误差失控，不得外推。解析名称只有在方程、参数和边界真正匹配时才成立。

目标本身还可能不可同时满足。若事实准确、安全边界和响应时限形成空可行集，任何势的下降都无法产生合法终点。控制器应先做可行性检查，无法满足时返回冲突证据并请求目标修订，而不是增大权重掩盖矛盾。多目标权重法只能找某个标量化问题的解，非凸情况下还可能漏掉 Pareto 前沿；因此关键约束应显式保留，软目标才进入势函数。对随时间变化的偏好，记录谁在何时改动了权重以及对已有承诺的影响。

随机策略下，期望势下降也不保证每条轨迹安全。应分别给出期望性能、尾部风险和边界触碰概率，必要时采用机会约束或屏障函数。若外部行动不可逆，仿真中的回滚不能抵消真实后果，提交前必须经过影子执行或权限门。AgentOS 的 TCE 可以提出目标更新，但 Boundary 与环境回执拥有独立否决权；这体现目标生成和行动资格的分离，也防止“目的”在模型中被误写成无条件外力。

\section{可解析基准的验证边界}
\label{sec:conclusion_first_analytical_solution}

本章得到的可靠结果可以分层陈述。定义层包括伯努利状态、Fisher 度量、代价单位、算子域和读出接口；模型假设包括固定命题身份、一维充分统计量、给定势与边界；数学结果包括等距坐标、欧拉—拉格朗日方程以及特定边界下的谱；经验主张则需要数据验证，例如有效惯性是否稳定、编辑延迟是否受图路径控制；工程设计包括 AgentOS 的模块分工、日志与回归门槛。把这些层混在一起，会由“方程可解”跳到“世界如此”，正是本章要避免的错误。

最小验证套件应包含四类测试。第一，符号测试：验证 Fisher 度量、坐标变换、变分方程和参数量纲。第二，数值测试：网格加密、本征值收敛、连续与离散轨迹对齐，并扫描步长稳定区。第三，机制测试：分别改变结构图延迟、读出器、目标和输入，确认模型能区分导致相同 $q(t)$ 的不同原因。第四，任务测试：同时报告准确率、校准、约束违规、资源代价和无关能力回归。任一项失败都不能被另一项的漂亮曲线抵消。

还应设置明确的反例。高置信错误表明低熵不等于正确；Dirichlet、Neumann 与周期边界的不同谱表明离散模不是唯一语义实体；快速广播的大图表明拓扑规模不等于时间；目标势下降而任务失败表明优化代理不等于成功；文本读出正确而结构冲突表明全局分布表示不能替代局部结构审计。这些反例不是削弱 HSF-HD，而是给一般状态动力学划定可证范围。理论作为智能的一般“空气动力学”，应提供跨实现的变量关系、稳定性问题和验证方法，而不是把一个 AgentOS 架构写成所有智能的唯一机制。

复现实验时，数据包至少保存命题定义、模型和提示版本、随机种子、观测时间戳、边界类型、单位表、离散步长、初值、目标版本、结构快照以及读出代码。训练数据不可得时，应明确把结论限制为行为辨识。若外部记忆或工具会变化，需固定快照或记录响应。报告中将解析解与数值误差分开，将统计置信区间与状态区间分开，将代价单位与焦耳电耗分开。这样别人才能判断差异来自模型、实现还是环境。

伯努利模型的适用边界同样清楚：它不表示多命题相关性、组合身份、层级计划、工具交互和长期记忆竞争；端点附近的有限样本估计敏感；真实系统的目标、结构与度量往往随时间变化。超出这些条件时，应升级到多项分布、指数族、图耦合状态或混合离散—连续系统，而不是继续给一维变量附加比喻。升级后仍可保留本章的纪律：先定义对象与单位，再写假设与方程，随后推导并用反例测试。

因此，这个“第一可解基准”的地位不是本体论确证，而是校准尺。它让我们在最简单空间中看清坐标奇性、量纲闭合、边界依赖、算子域、目标迁移和离散实现，再把这些检查带回高维模型。只要它能持续发现错误、约束接口并给出可重复预测，它就完成了任务；若未来证据否定某个机制，我们应修改机制而不是维护隐喻。HSF-HD 的力量来自可修正的计算动力学框架，AgentOS 的价值来自可审计的工程实例，二者相关而不等同。
